\documentclass[11pt]{article}
\usepackage[margin=0.82in]{geometry}
\usepackage{amsmath,amssymb,lmodern}
\usepackage[T1]{fontenc}
\usepackage{needspace}
\usepackage[colorlinks=true,linkcolor=blue,urlcolor=blue]{hyperref}
\setlength{\emergencystretch}{3em}
\setlength{\parskip}{0.3em}
\setcounter{secnumdepth}{0}
\providecommand{\tightlist}{\setlength{\itemsep}{0pt}\setlength{\parskip}{0pt}}
\title{Heat-content maximization on compact metric graphs at every time}
\author{Henry Zweiman}
\date{September 16, 2026}
\begin{document}
\maketitle
\Needspace{8\baselineskip}\section{Abstract}\label{abstract}

Among finite compact connected metric graphs of prescribed total length
with a nonempty Dirichlet vertex set and standard conditions elsewhere,
the interval with one Dirichlet and one Neumann endpoint maximizes heat
content at every positive time. Equality at a single positive time
characterizes this interval, up to subdivision. This answers the
all-times question of Bifulco and Täufer. We prove the stronger
statement that heat evolution on any such graph is dominated in
cumulative decreasing rearrangement by mixed-boundary interval evolution
of the rearranged initial datum. The argument combines the classical
metric-graph level-set inequality with the zero-order elliptic
concentration comparison and implicit-time iteration of parabolic
symmetrization. The rigidity proof uses the half-time energy identity
and strict edgewise concavity of the heat evolution of the constant
function. No arithmetic condition on the edge lengths is imposed. We
also prove that, among graphs of fixed length carrying a finite
nonnegative killing measure of fixed positive mass, the interval with
all killing at one endpoint maximizes heat content at every time. This
comparison allows singular measures and arbitrary nonnegative initial
data, and equality at a single positive time characterizes the interval
and its endpoint measure. For the same general killing measures we prove
sharp Kohler-Jobin inequalities at fixed length and total killing, and
at fixed torsion and total killing. A ground-state distribution estimate
controls the potential terms and compares the resulting scalar
differential inequality with the cosine profile; both equality cases
characterize the endpoint-killed interval. The stationary joint
comparisons also extend to the nonlinear \(p\)-Laplacian for every
\(1<p<\infty\), with explicit finite-strength constants and the same
rigidity.

\textbf{Keywords:} metric graph; heat content; rearrangement;
concentration comparison; Faber-Krahn inequality; rigidity.

\textbf{Mathematics Subject Classification (2020):} 35K05, 34B45, 49Q10.

\Needspace{8\baselineskip}\section{1. Introduction and main
results}\label{introduction-and-main-results}

Let \(\Gamma\) be a finite compact connected metric graph, with positive
edge lengths and total length \(L>0\). Let \(D\) be a nonempty subset of
its vertices. We assume that \(\Gamma\setminus D\) is connected. The
Laplacian has Dirichlet conditions at \(D\) and standard continuity and
Kirchhoff conditions at every other vertex. A standard condition at a
vertex of degree one is a Neumann condition. Loops and multiple edges
are permitted; degrees and Kirchhoff sums count edge incidences, so a
loop contributes twice.

Write \(A_\Gamma\) for the nonnegative realization of \(-d^2/dx^2\) and
\(S_\Gamma(t)=e^{-tA_\Gamma}\). Its heat content is

\[
Q_\Gamma(t)=\int_\Gamma S_\Gamma(t)1\,dx.
\tag{1}
\]

We compare \(\Gamma\) to \(I_L=(0,L)\), with Neumann condition at \(0\)
and Dirichlet condition at \(L\). Denote the corresponding operator,
semigroup and heat content by \(A_I\), \(S_I(t)\) and \(Q_I(t)\).

\Needspace{6\baselineskip}\textbf{Theorem 1.1 (heat-content maximization and rigidity).} For every
graph described above and every \(t>0\),

\[
Q_\Gamma(t)\le Q_I(t).
\tag{2}
\]

Equality for one positive time holds if and only if \(\Gamma\) is a path
of length \(L\), \(D\) consists of one endpoint, and all other vertices
have standard conditions. Equivalently, after suppressing vertices of
degree two, the graph is the Dirichlet-Neumann interval. In particular,
for every other admissible graph the inequality is strict at every
positive time.

For a nonnegative measurable function \(f\) on a space of length \(L\),
let \(f^*\) be its decreasing rearrangement on \((0,L)\), and set

\[
\mathcal C_f(s)=\int_0^s f^*(r)\,dr,
\qquad 0\le s\le L.
\tag{3}
\]

We write \(f\preccurlyeq g\) if \(\mathcal C_f(s)\le\mathcal C_g(s)\)
for all \(s\). Equality of total integrals is not part of this
definition. Thus the order is weak submajorization, also called
concentration order in the parabolic comparison literature.

\Needspace{6\baselineskip}\textbf{Theorem 1.2 (concentration comparison).} Let
\(0\le f\in L^2(\Gamma)\) and let \(0\le g\in L^2(0,L)\) be
nonincreasing. If

\[
\mathcal C_f(s)\le\int_0^s g(r)\,dr
\quad\text{for every }s\in[0,L],
\tag{4}
\]

then, for every \(t>0\),

\[
\mathcal C_{S_\Gamma(t)f}(s)
\le\int_0^s S_I(t)g(r)\,dr
\quad\text{for every }s\in[0,L].
\tag{5}
\]

The interval function on the right is nonnegative and nonincreasing. In
particular, one may take \(g=f^*\).

Taking \(f=g=1\) and \(s=L\) proves the inequality in Theorem 1.1. Its
one-time equality assertion needs an additional argument and is proved
in Section 5.

\subsection{1.1. Relation to earlier
work}\label{relation-to-earlier-work}

Bifulco and Täufer {[}BT, Theorem 2.6{]} prove the interval comparison
for sufficiently large times, and for sufficiently small times when the
edge lengths are rationally dependent. Their final Remark 2.9 leaves the
comparison at every positive time open. Theorem 1.1 resolves that
question without a restriction on edge lengths or on the time. The
graph, vertex conditions and total-length normalization here agree with
their Definitions 2.1 and 2.2. Their comparison interval is written with
its Dirichlet endpoint at \(0\); reflection gives our convention.

Bifulco and Mugnolo {[}BM{]} derive a heat-content formula, short-time
asymptotics and surgery principles for compact quantum graphs. Their
Section 6 includes comparisons under imposing Dirichlet conditions,
taking subgraphs and attaching pendant graphs, as well as
heat-content-preserving constructions. These are a different set of
comparisons from (5), which rearranges arbitrary initial data and
compares graphs of the same total length. We use the ordinary Kirchhoff
heat semigroup studied in that work.

The geometric rearrangement input is classical. In the proof of {[}F,
Lemma 3{]}, Friedlander uses the number of points in a regular level set
to compare the graph Dirichlet energy to the energy of a monotone
rearrangement on an interval. That proof also explains how equality in
level multiplicity excludes branching. We give the relevant calculation
because the zero-order term of the resolvent equation must be retained.

The passage from elliptic concentration comparison to parabolic
comparison by implicit-time discretization is likewise classical.
Vázquez {[}V, Sections 5 and 7{]} develops precisely this strategy for
Euclidean equations; see in particular the zero-order comparison and the
parabolic iteration in the author\textquotesingle s available text. We
adapt that mechanism to the mixed-boundary graph problem. Neither the
level-set Cauchy-Schwarz inequality nor implicit Euler iteration is
asserted to be new. The contribution is the all-times graph comparison
(5), its answer to {[}BT, Remark 2.9{]}, and the one-time rigidity
theorem. The elementary consequences in Section 6 are included as
applications of the same result. Section 7 extends the concentration
theorem to nonnegative killing measures of prescribed total mass, using
a Robin interval comparator and a boundary inequality for cumulative
resolvents. It includes delta-vertex conditions, nonnegative potentials
and finite singular measures.

\Needspace{8\baselineskip}\section{2. Graph forms and rearrangement
preliminaries}\label{graph-forms-and-rearrangement-preliminaries}

The measure on \(\Gamma\) is the sum of Lebesgue measures on the edges;
vertices have measure zero. The closed quadratic form of \(A_\Gamma\) is

\[
\mathfrak a_\Gamma[u]=\sum_{e}\int_0^{\ell_e}|u_e'(x)|^2\,dx,
\qquad
\mathcal D(\mathfrak a_\Gamma)
=\{u\in H^1(\Gamma):u|_D=0\}.
\tag{6}
\]

Here \(H^1(\Gamma)\) consists of edgewise \(H^1\) functions continuous
at the vertices. The operator domain consists of edgewise \(H^2\)
functions in the form domain satisfying the Kirchhoff sum at each vertex
outside \(D\). The same construction on \(I_L\), with form domain
\(\{v\in H^1(0,L):v(L)=0\}\), gives \(A_I\).

Both operators are nonnegative and self-adjoint, with compact resolvent.
They are strictly positive: if \(p\in D\), integration along a simple
path from \(p\) to any point gives

\[
|u(x)|^2\le L\,\mathfrak a_\Gamma[u],
\qquad
\|u\|_2^2\le L^2\mathfrak a_\Gamma[u].
\tag{7}
\]

The truncations \(u\mapsto u_+\) and \(u\mapsto\min\{1,\max\{u,0\}\}\)
preserve the form domain and do not increase its energy. The resolvents
and semigroups are therefore positivity preserving and sub-Markovian. In
particular,

\[
0\le S_\Gamma(t)1\le1,
\qquad
0\le S_I(t)1\le1.
\tag{8}
\]

We recall two elementary rearrangement facts. For nonnegative
\(f\in L^1\) on a nonatomic space of length \(L\),

\[
\mathcal C_f(s)=\sup_{|E|=s}\int_E f,
\qquad
\sup_{0\le s\le L}|\mathcal C_f(s)-\mathcal C_h(s)|
\le\|f-h\|_1.
\tag{9}
\]

The supremum formula follows by filling a superlevel set of \(f\) and,
if necessary, part of one level set. The second assertion follows by
comparing the integrals on each candidate set. Also,

\[
\|f^*-h^*\|_{L^2(0,L)}\le\|f-h\|_{L^2(\Gamma)}
\tag{10}
\]

for nonnegative \(f,h\). Indeed, Hardy-Littlewood gives
\(\int_\Gamma fh\le\int_0^L f^*h^*\), while rearrangement preserves the
two squared norms. Expanding the squared distances proves (10).

\Needspace{6\baselineskip}\textbf{Lemma 2.1 (a Lipschitz rearrangement).} Suppose \(u\ge0\) is
continuous on \(\Gamma\), vanishes at \(D\), and is Lipschitz on every
edge, with maximum edge Lipschitz constant \(K\). Then \(u^*\) has a
Lipschitz representative on \([0,L]\) with Lipschitz constant at most
\(K\), and \(u^*(L)=0\).

\textbf{Proof.} The case \(K=0\) is immediate. Choose a maximum point
and a simple path from it to \(D\). For \(0\le c_1<c_2\le\max u\), this
path contains a subarc along which \(u\) passes from \(c_2\) to \(c_1\)
while remaining between those values. The length of that subarc is at
least \((c_2-c_1)/K\). The path is simple, so its length is counted
without multiplicity in the graph measure. Consequently

\[
|\{c_1<u<c_2\}|\ge\frac{c_2-c_1}{K}.
\tag{11}
\]

This inverse-distribution estimate gives the Lipschitz assertion; jumps
of the distribution due to plateaus correspond to constant portions of
the rearrangement and cause no difficulty. Finally, continuity at a
Dirichlet vertex shows that the essential infimum is zero. Hence the
representative has \(u^*(L)=0\). \(\square\)

\Needspace{8\baselineskip}\section{3. Resolvent concentration
comparison}\label{resolvent-concentration-comparison}

For \(h>0\), put

\[
R_h=(1+hA_\Gamma)^{-1},
\qquad B_h=(1+hA_I)^{-1}.
\tag{12}
\]

\Needspace{6\baselineskip}\textbf{Lemma 3.1 (the interval resolvent).} If \(g\ge0\) is
nonincreasing and belongs to \(L^2(0,L)\), then \(v=B_hg\) is
nonnegative and nonincreasing. With \(V(s)=\int_0^s v\) and
\(G(s)=\int_0^s g\),

\[
-hV''+V=G,
\qquad V(0)=0,\qquad V'(L)=0.
\tag{13}
\]

If \(g_1,g_2\) are two such functions and
\(\int_0^s g_1\le\int_0^s g_2\) for all \(s\), then
\(\int_0^s B_hg_1\le\int_0^s B_hg_2\) for all \(s\).

\textbf{Proof.} First let \(g\) be smooth on \([0,L]\). The equation is
\(-hv''+v=g\), with \(v'(0)=0\) and \(v(L)=0\). Positivity of the
resolvent gives \(v\ge0\), hence \(v'(L)\le0\). Its derivative \(q=v'\)
satisfies

\[
-hq''+q=g'\le0,\qquad q(0)=0,\quad q(L)\le0.
\tag{14}
\]

Testing against \(q_+\) gives \(q\le0\). Every nonnegative nonincreasing
\(L^2\) function can be approximated in \(L^2\) by smooth nonnegative
nonincreasing functions: first truncate, then extend by constant
endpoint values and mollify. Boundedness of \(B_h\) on \(L^2\) passes
the conclusion to the limit; the cone of nonincreasing functions is
closed in \(L^2\).

Integration of the resolvent equation from \(0\) to \(s\), using
\(v'(0)=0\), proves (13). This remains valid for \(L^2\) data because
\(v\in H^2\). For the last assertion, the difference \(W=V_1-V_2\)
satisfies \(-hW''+W\le0\), \(W(0)=0\) and \(W'(L)=0\). Thus

\[
h\int_0^L|(W_+)'|^2+\int_0^L|W_+|^2\le0,
\tag{15}
\]

and \(W\le0\). \(\square\)

\Needspace{6\baselineskip}\textbf{Lemma 3.2 (graph resolvent comparison).} For every
\(0\le f\in L^2(\Gamma)\),

\[
\mathcal C_{R_hf}(s)\le\int_0^s B_h(f^*)(r)\,dr
\quad(0\le s\le L).
\tag{16}
\]

Consequently, if \(g\) satisfies the hypotheses and concentration bound
of Theorem 1.2, then

\[
\mathcal C_{R_hf}(s)\le\int_0^s B_hg(r)\,dr.
\tag{17}
\]

\textbf{Proof.} We first suppose that on every closed edge the load
\(f\) is a strictly positive nonconstant polynomial. Loads need not be
continuous across vertices. The solution \(u=R_hf\) satisfies

\[
-hu_e''+u_e=f_e
\tag{18}
\]

on each edge, so it extends analytically past both endpoints as an
edgewise function. It is strictly positive outside \(D\). For
completeness, a zero interior minimum contradicts (18) and \(f_e>0\). At
a non-Dirichlet vertex with value zero, all outgoing derivatives would
be nonnegative, and their Kirchhoff sum would make them all zero.
Equation (18) would then give a negative second derivative into each
incident edge at that endpoint, contradicting nonnegativity. This proves
the positivity assertion.

The function \(u\) cannot be constant on an edge, because (18) would
force \(f\) to be constant there. Each edge therefore has only finitely
many critical points: its derivative is analytic on a neighborhood of
the closed edge and is not identically zero. Exclude their values and
all vertex values. For every remaining \(c\in(0,\max u)\), write

\[
E_c=\{u>c\},\quad \mu(c)=|E_c|,\quad
N(c)=\#\{u=c\},
\quad a(c)=\sum_{u=c}|u'|,\quad
b(c)=\sum_{u=c}\frac1{|u'|}.
\tag{19}
\]

All sums run over interior edge points. Continuity along a path from a
maximum to \(D\) gives \(N(c)\ge1\). Define \(U=\mathcal C_u\) and
\(F=\mathcal C_f\). Integration of (18) over \(E_c\) gives

\[
ha(c)+U(\mu(c))=\int_{E_c}f\le F(\mu(c)).
\tag{20}
\]

To check the sign and vertex terms, the derivative pointing out of a
superlevel interval is \(-|u'|\) at a regular boundary point. At a
vertex contained in \(E_c\), continuity puts all incident edge germs in
\(E_c\), and their fluxes cancel by Kirchhoff. No Dirichlet vertex
belongs to \(E_c\). These observations give exactly the positive flux
term on the left of (20).

The one-dimensional coarea formula and inverse differentiation give

\[
-\mu'(c)=b(c),\qquad
U''(\mu(c))=(u^*)'(\mu(c))=-\frac1{b(c)}.
\tag{21}
\]

Cauchy-Schwarz, in the form already used for graph rearrangement in
{[}F{]}, yields

\[
a(c)b(c)\ge N(c)^2\ge1.
\tag{22}
\]

Combining (20)-(22), we obtain

\[
-hU''+U\le F.
\tag{23}
\]

There is no omitted singular part in (23). The function \(u^*\) is
Lipschitz by Lemma 2.1, so \(U\in W^{2,\infty}(0,L)\). There are no
plateaus of \(u\), and hence the finitely many exceptional levels yield
only finitely many exceptional ranks \(\mu(c)\). Thus the displayed
inequality holds almost everywhere and in the weak sense. Its boundary
data are

\[
U(0)=0,\qquad U'(L)=0.
\tag{24}
\]

By Lemma 3.1, the cumulative integral \(V\) of \(B_h(f^*)\) satisfies
the equality in (23), with the same boundary data. The positive-part
test (15) proves \(U\le V\).

To remove the polynomial assumption, approximate an arbitrary
nonnegative \(L^2\) load by strictly positive nonconstant polynomials on
every closed edge. Such loads are dense: separately on each edge,
approximate first by a nonnegative continuous function, then uniformly
by a polynomial, adding a vanishing positive constant to ensure strict
positivity. A vanishing linear perturbation makes a constant polynomial
nonconstant while preserving positivity. There are only finitely many
edges.

For an approximating sequence \(f_n\to f\) in \(L^2\), resolvent
boundedness and (10) give

\[
R_hf_n\to R_hf,\qquad
B_h(f_n^*)\to B_h(f^*)
\quad\text{in }L^2.
\tag{25}
\]

Finite total length converts these to \(L^1\) convergences, and (9)
passes the cumulative inequality to the limit. This proves (16). Finally
apply Lemma 3.1 to the decreasing interval data \(f^*\) and \(g\) to
obtain (17). \(\square\)

The comparison retains the zero-order term in (18). Dropping that term
would give a Poisson comparison, which does not by itself iterate to the
required heat inequality.

\Needspace{8\baselineskip}\section{4. Passage to the heat
semigroup}\label{passage-to-the-heat-semigroup}

\textbf{Proof of Theorem 1.2.} Fix \(h>0\) and define

\[
u_0=f,\quad u_j=R_hu_{j-1},\qquad
v_0=g,\quad v_j=B_hv_{j-1}.
\tag{26}
\]

All these functions are nonnegative, and every \(v_j\) is nonincreasing
by Lemma 3.1. The hypothesis at \(j=0\) and (17) give inductively

\[
\mathcal C_{u_j}(s)\le\int_0^s v_j(r)\,dr
\quad(0\le s\le L,\ j\ge0).
\tag{27}
\]

For a fixed \(t>0\), take \(h=t/n\) and \(j=n\). The spectral theorem
gives

\[
\left(1+\frac tn A_\Gamma\right)^{-n}f\longrightarrow S_\Gamma(t)f,
\qquad
\left(1+\frac tn A_I\right)^{-n}g\longrightarrow S_I(t)g
\quad\text{in }L^2.
\tag{28}
\]

Indeed, \((1+t\lambda/n)^{-n}\to e^{-t\lambda}\) for every
\(\lambda\ge0\), and both multipliers are bounded by one. Dominated
convergence in the spectral measure proves (28). Passing to the limit in
(27) by (9) proves (5), uniformly in \(s\). Closedness of the cone of
nonincreasing functions gives the remaining monotonicity assertion.
\(\square\)

This is the classical implicit-time strategy of {[}V{]}, with the graph
resolvent inequality supplied by Section 3. In particular, no
differentiability of a time-dependent rearrangement and no assumption
about regular heat levels are needed to prove the inequality.

\Needspace{8\baselineskip}\section{5. Equality at a single time}\label{equality-at-a-single-time}

The proof of rigidity uses the constant initial datum. Put

\[
w(t)=S_\Gamma(t)1,\qquad z(t)=S_I(t)1.
\tag{29}
\]

We first record a strict regularity property of this particular heat
evolution.

\Needspace{6\baselineskip}\textbf{Lemma 5.1 (strict edgewise concavity).} For every \(t>0\), the
function \(w(t)\) is smooth on the closure of each individual edge,
positive outside \(D\), and

\[
-w_e''(t,x)=A_\Gamma w(t,x)>0
\quad(0<x<\ell_e).
\tag{30}
\]

Consequently, each open edge has at most one critical point of \(w(t)\)
and has no constant subinterval.

\textbf{Proof.} Spectral smoothing gives
\(w(t)\in\mathcal D(A_\Gamma^m)\) for every integer \(m\ge1\). Repeated
use of the edge equation for the operator gives edgewise \(H^{2m}\)
regularity, and hence smoothness up to each endpoint. On a graph
connected outside its Dirichlet set, the Kirchhoff heat semigroup is
positivity improving on \(\Gamma\setminus D\).

Here is a maximum-principle justification of the positivity statement
used in this proof. For nonzero nonnegative initial data, positivity
preservation and injectivity of \(S_\Gamma(t)\) give a nonzero
nonnegative profile at positive times. An interior zero at a positive
time forces the solution to vanish backwards on that edge by the
parabolic strong maximum principle. At a standard vertex where the
nonnegative solution has value zero, all outgoing derivatives are
nonnegative. Their sum is zero, so every one vanishes; the parabolic
boundary point lemma excludes a positive incident edge with zero
outgoing derivative. Thus positivity propagates across each standard
vertex. Connectivity of \(\Gamma\setminus D\) gives positivity
throughout the open edges and at its standard vertices. This is the
usual strict-positivity property of the connected graph heat kernel; see
also the discussion in {[}BM, Section 1{]}.

Sub-Markovianity and positivity give

\[
w(t+h)=S_\Gamma(t)S_\Gamma(h)1\le S_\Gamma(t)1=w(t)
\quad(h>0).
\tag{31}
\]

Differentiation at positive time in \(L^2\) implies
\(A_\Gamma w(t)=-\partial_t w(t)\ge0\). It is not the zero function:
otherwise (7) would give \(w(t)=0\), which contradicts positivity.
Commutation through the semigroup yields

\[
A_\Gamma w(t)=S_\Gamma(t/2)A_\Gamma w(t/2).
\tag{32}
\]

Positivity improvement applied to the nonzero nonnegative datum on the
right proves (30). The final assertion follows because \(w_e'\) is
strictly decreasing. \(\square\)

\Needspace{6\baselineskip}\textbf{Lemma 5.2 (equal norms under concentration).} Suppose
\(a,b\ge0\) belong to \(L^2\) on spaces of the same finite measure,
\(a\preccurlyeq b\), and \(\|a\|_2=\|b\|_2\). Then \(a^*=b^*\) almost
everywhere.

\textbf{Proof.} For \(c\ge0\) define

\[
H_a(c)=\int(a-c)_+
=\sup_{0\le s\le L}\{\mathcal C_a(s)-cs\}.
\tag{33}
\]

The same formula holds for \(b\), so \(H_a\le H_b\). Tonelli gives

\[
\|a\|_2^2=2\int_0^\infty H_a(c)\,dc,
\qquad
\|b\|_2^2=2\int_0^\infty H_b(c)\,dc.
\tag{34}
\]

Equality of the norms makes \(H_a=H_b\) almost everywhere, and then
everywhere by continuity. Their right derivatives are \(-|\{a>c\}|\) and
\(-|\{b>c\}|\), respectively. Thus the distributions agree and so do the
rearrangements. \(\square\)

\textbf{Proof of the rigidity assertion in Theorem 1.1.}
Self-adjointness and the semigroup law give the half-time identities

\[
Q_\Gamma(t)=\|w(t/2)\|_2^2,
\qquad
Q_\Gamma'(t)=-\mathfrak a_\Gamma[w(t/2)].
\tag{35}
\]

The identical formulas hold on the interval. Suppose
\(Q_\Gamma(T)=Q_I(T)\) for some \(T>0\). Theorem 1.2 at time \(T/2\)
gives \(w(T/2)\preccurlyeq z(T/2)\). Lemma 5.2 and (35) imply

\[
w(T/2)^*=z(T/2)
\quad\text{almost everywhere on }(0,L).
\tag{36}
\]

The differentiable nonnegative function \(Q_I(t)-Q_\Gamma(t)\) attains
its minimum zero at the interior point \(T\), so its derivative vanishes
there. Thus (35) gives

\[
\mathfrak a_\Gamma[w(T/2)]
=\mathfrak a_I[z(T/2)]
=\int_0^L|(w(T/2)^*)'|^2.
\tag{37}
\]

Set \(u=w(T/2)\) and \(M=\max u\). By Lemma 5.1, only finitely many
critical or vertex values need be excluded. At all other levels
\(c\in(0,M)\) use the notation \(a(c),b(c),N(c)\) from (19). Lemma 2.1
and one-dimensional coarea yield

\[
\mathfrak a_\Gamma[u]=\int_0^M a(c)\,dc,
\qquad
\int_0^L|(u^*)'|^2=\int_0^M\frac1{b(c)}\,dc.
\tag{38}
\]

The second formula follows from \(s=\mu(c)\), \(ds=-b(c)dc\) and
\((u^*)'(\mu(c))=-1/b(c)\) on every interval between exceptional values.
The Lipschitz representative ensures that these identities lose no
singular energy. By (22),

\[
a(c)-\frac1{b(c)}\ge\frac{N(c)^2-1}{b(c)}\ge0.
\tag{39}
\]

Equations (37)-(39) force \(N(c)=1\) for almost every regular \(c\).

Every edge incidence at a Dirichlet vertex has strictly positive
derivative into its edge. Indeed, the function is positive inside that
edge, zero at its endpoint, and strictly concave; the endpoint
derivative is at least any positive chord slope. Therefore, for all
sufficiently small positive regular levels,

\[
N(c)\ge\sum_{p\in D}\deg(p).
\tag{40}
\]

The small incident edge neighborhoods can be chosen disjoint, including
the two incidences of a loop. It follows that \(D\) consists of one
vertex and that this vertex has degree one.

Next consider a standard vertex \(p\) of degree \(d\ge3\). Write
\(c=u(p)>0\) and let \(d_i\) be the derivatives along its outgoing
incidences. Kirchhoff gives \(\sum_i d_i=0\). Each incidence with
\(d_i>0\) gives increasing values immediately away from \(p\); each
incidence with \(d_i<0\) gives decreasing values. If \(d_i=0\), strict
edgewise concavity also gives decreasing values. Thus every incidence is
assigned to one of two classes. Since \(d\ge3\), one class contains at
least two incidences. On the corresponding side of \(c\), all
sufficiently close regular levels have at least two distinct preimages.
This contradicts \(N=1\) almost everywhere. No standard vertex can have
degree at least three.

A connected finite graph whose degrees are at most two and which has a
vertex of degree one is a path. Its unique Dirichlet vertex is an
endpoint; its other endpoint is standard and hence Neumann. Suppressing
standard degree-two vertices changes neither the metric space nor its
Laplacian. This proves necessity. Sufficiency follows by reflection and
subdivision invariance of the interval problem. \(\square\)

The equality argument does not infer strictness from strict resolvent
inequalities followed by a limit. Such a limit could erase a positive
defect. Instead, equality at the specified time produces the exact
energy equality (37).

\Needspace{8\baselineskip}\section{6. Consequences and limits}\label{consequences-and-limits}

\Needspace{6\baselineskip}\textbf{Corollary 6.1 (convex temperature functionals).} Under the
hypotheses of Theorem 1.2, for every convex nondecreasing function
\(\Phi:[0,\infty)\to[0,\infty)\) with \(\Phi(0)=0\),

\[
\int_\Gamma\Phi(S_\Gamma(t)f)\,dx
\le\int_0^L\Phi(S_I(t)g)\,dx.
\tag{41}
\]

In particular, all \(L^p\) norms, \(1\le p\le\infty\), are bounded by
those of the comparison evolution.

\textbf{Proof.} Concentration order implies the stop-loss comparison
(33) and the comparison of first moments. Every such convex \(\Phi\) is
a nonnegative linear function plus an integral of \((r-c)_+\) against
its distributional second derivative; approximation applies if
necessary. Integrating that representation proves (41). Choosing
\(\Phi(r)=r^p\) gives finite \(p\); passage to the limit gives
\(p=\infty\). At positive time the profiles are bounded by spectral
smoothing and the edge Sobolev embedding. \(\square\)

The interval bound has the explicit expansion

\[
Q_I(t)=\frac{8L}{\pi^2}\sum_{k=0}^{\infty}
\frac{\exp\!\left(-\frac{(2k+1)^2\pi^2t}{4L^2}\right)}{(2k+1)^2}.
\tag{42}
\]

Indeed, the normalized eigenfunctions are
\(\sqrt{2/L}\cos((2k+1)\pi x/(2L))\). Their squared integrals are
\(8L/(\pi^2(2k+1)^2)\), giving (42) by the spectral representation of
(1).

\Needspace{6\baselineskip}\textbf{Corollary 6.2 (positive time weights).} Let
\(\rho:(0,\infty)\to[0,\infty)\) be measurable. Whenever the interval
integral is finite,

\[
\int_0^\infty\rho(t)Q_\Gamma(t)\,dt
\le\int_0^\infty\rho(t)Q_I(t)\,dt.
\tag{43}
\]

If \(\rho\) is positive on a set of positive measure, equality
characterizes the same interval. In particular, choosing \(\rho=1\)
recovers the sharp torsional-rigidity bound
\(\int_0^\infty Q_\Gamma(t)dt\le L^3/3\).

\textbf{Proof.} Integrate Theorem 1.1 and use its strictness. For the
last constant, the interval torsion function solves \(-v''=1\),
\(v'(0)=v(L)=0\), hence \(v(x)=(L^2-x^2)/2\) and \(\int_0^L v=L^3/3\).
\(\square\)

The torsional comparison itself is prior work; it is discussed in {[}BT,
Remark 2.3{]} and is not a separate novelty claim here. Similarly, with
the diffusion convention that the generator is \(d^2/dx^2\) on each
edge, \(Q_\Gamma(t)/L\) is the survival probability at time \(t\) for
the killed graph diffusion started from uniform length measure. Theorem
1.1 therefore also says that its lifetime is stochastically dominated by
the lifetime on the mixed interval. This is the direct probabilistic
interpretation of the heat-content bound, using the Feynman-Kac
representation in {[}BT, Proposition 3.2{]}.

The conclusions of Sections 1-6 concern unweighted finite metric graphs
with the stated Dirichlet and Kirchhoff conditions. Section 7 gives a
separate theorem for nonnegative killing measures with continuous form
domain. They do not provide a heat-trace bound, pointwise heat-kernel
domination, or a theorem for arbitrary vertex couplings. We prove no
quantitative stability estimate. A useful next problem is to quantify
the loss in (2) in terms of geometric distance from a path, with
explicit control of the time scale and of edges whose lengths tend to
zero. Such a statement requires additional work beyond the qualitative
rigidity proved here.

\Needspace{8\baselineskip}\section{7. A fixed amount of nonnegative
killing}\label{a-fixed-amount-of-nonnegative-killing}

The comparison extends to a different boundary regime: there are no
Dirichlet conditions, and killing is described by a finite nonnegative
measure. Throughout this section, \(\Gamma\) is a finite compact
connected metric graph of total length \(L\), and \(\nu\) is a
nonnegative Radon measure on \(\Gamma\) with

\[
\nu(\Gamma)=\kappa>0.
\tag{44}
\]

The measure \(\nu\) is distinct from the length measure \(dx\) used for
heat content and rearrangement. In particular, \(\nu\) may charge
vertices. On the space \(H^1(\Gamma)\) of continuous edgewise \(H^1\)
functions, define

\[
\mathfrak a_\nu[u]=\int_\Gamma |u'|^2\,dx+\int_\Gamma |u|^2\,d\nu,
\qquad S_\nu(t)=e^{-tA_\nu},
\qquad Q_\nu(t)=\langle 1,S_\nu(t)1\rangle.
\tag{45}
\]

Here \(A_\nu\) is the operator associated with the form. This includes
nonnegative integrable potentials and nonnegative delta interactions,
separately or together, as well as singular continuous killing measures.
It does not include signed potentials or weighted continuity conditions.

The comparison form on \(I_L=(0,L)\) is

\[
\mathfrak a_{I,\kappa}[v]=\int_0^L|v'|^2\,dx+\kappa|v(L)|^2,
\qquad \mathcal D(\mathfrak a_{I,\kappa})=H^1(0,L).
\tag{46}
\]

Its operator has \(v'(0)=0\) and \(v'(L)+\kappa v(L)=0\). Write
\(S_{I,\kappa}\) and \(Q_{I,\kappa}\) for its semigroup and heat
content.

\Needspace{6\baselineskip}\textbf{Theorem 7.1 (fixed-total-killing comparison and rigidity).}
Suppose \(\nu\) satisfies (44), \(0\le f\in L^2(\Gamma)\), and
\(0\le g\in L^2(0,L)\) is nonincreasing with \(f\preccurlyeq g\). Then
\(S_{I,\kappa}(t)g\) is nonnegative and nonincreasing, and

\[
\mathcal C_{S_\nu(t)f}(s)
\le \int_0^s S_{I,\kappa}(t)g(r)\,dr
\quad(0\le s\le L,\ t>0).
\tag{47}
\]

In particular,

\[
Q_\nu(t)\le Q_{I,\kappa}(t)\qquad(t>0).
\tag{48}
\]

Equality in (48) at one positive time holds if and only if \(\Gamma\) is
a path of length \(L\) and \(\nu=\kappa\delta_p\) at one endpoint \(p\).
Subdivision and reflection are immaterial. Thus the interval with all
killing concentrated at one endpoint is the unique maximizer, as a
metric space with killing measure, among the stated graphs and measures
with fixed \(L\) and \(\kappa\).

For killing supported on finitely many vertices, the time integral of
(48) is the already known sharp torsion inequality of Özcan and Täufer
{[}OT, Theorem 8.2 in the cited preprint{]}. Their proof already places
all strength at a minimum and uses energy rearrangement. Our additional
assertion is the all-times concentration comparison (47), its extension
to arbitrary finite measures, and rigidity from equality at one time.
The key extra ingredient is a boundary inequality for the cumulative
resolvent, rather than an energy comparison alone. Generalized couplings
are also discussed in {[}BM, Remark 4.11{]}; that discussion does not
supply (47). The fixed-graph ground-state maximization problem of
Kurasov and Serio {[}KS{]} concerns a different functional and
optimization direction.

\subsection{7.1. Form and interval
preliminaries}\label{form-and-interval-preliminaries}

The compact embedding \(H^1(\Gamma)\hookrightarrow C(\Gamma)\) makes the
measure term in (45) continuous on \(H^1\). The form norm after adding
\(\|u\|_2^2\) is equivalent to the \(H^1\) norm, so the form is closed
and has compact resolvent. It is a Dirichlet form: real normal
contractions decrease both terms. Its semigroup and resolvents are
therefore positivity preserving and sub-Markovian.

The spectrum is strictly positive. Indeed, choose \(x_0\in\Gamma\) with
\(|u(x_0)|^2\le\kappa^{-1}\int|u|^2d\nu\), using continuity and the
minimum of \(|u|\). A simple path of length at most \(L\) from \(x_0\)
to \(x\) gives

\[
\|u\|_2^2\le 2L^2\int_\Gamma|u'|^2dx
+\frac{2L}{\kappa}\int_\Gamma|u|^2d\nu.
\tag{49}
\]

In this section Lemma 2.1 has the following version without Dirichlet
vertices: a continuous edgewise \(K\)-Lipschitz function has a
\(K\)-Lipschitz decreasing rearrangement with \(u^*(L)=\min_\Gamma u\).
Use a simple path from a maximum to a minimum in the proof of (11); only
levels between these two values are needed.

Fix \(h>0\) and let \(B_{h,\kappa}=(1+hA_{I,\kappa})^{-1}\). For
nonnegative nonincreasing \(g\), the function \(v=B_{h,\kappa}g\) is
nonnegative and nonincreasing. For smooth \(g\), its derivative
satisfies (14), now with \(q(L)=-\kappa v(L)\le0\); the same
positive-part argument applies. Monotone smooth approximation proves the
assertion for \(L^2\) data. For \(V(s)=\int_0^s v\) and
\(G(s)=\int_0^s g\) one has

\[
-hV''+V=G,\qquad V(0)=0,
\qquad V(L)+h\kappa V'(L)=G(L).
\tag{50}
\]

The differential equation follows by integration from the Neumann
endpoint, and its value at \(L\), together with \(v'(L)=-\kappa v(L)\),
gives the last condition. This is a Robin condition for \(V\) with an
inhomogeneous boundary value determined by the full integral of the
datum.

We shall use the following elementary comparison. If \(W\in H^2(0,L)\)
satisfies

\[
-hW''+W\le0,\qquad W(0)=0,
\qquad W(L)+h\kappa W'(L)\le0,
\tag{51}
\]

then integration against \(W_+\) gives

\[
h\int_0^L|(W_+)'|^2+\int_0^L W_+^2
+\frac{W_+(L)^2}{\kappa}\le0.
\tag{52}
\]

Thus \(W\le0\). In particular, (50) preserves cumulative order of
decreasing data: if \(G_1\le G_2\) on \([0,L]\), their solutions satisfy
\(V_1\le V_2\), including when \(G_1(L)<G_2(L)\).

\subsection{7.2. Resolvent comparison for finitely many killing
points}\label{resolvent-comparison-for-finitely-many-killing-points}

First suppose \(\nu=\sum_p\beta_p\delta_p\), with finitely many
\(\beta_p\ge0\) and \(\sum_p\beta_p=\kappa\). Insert degree-two vertices
at interior atoms. The operator acts as \(-u''\) on each open edge and
satisfies

\[
\sum_{e\sim p}\partial_{\mathrm{out},e}u(p)=\beta_pu(p).
\tag{53}
\]

Here each derivative points from the vertex into its incident edge; this
convention makes the derivative at the right endpoint of an interval
equal to \(-u'(L)\). Vertices without an atom have \(\beta_p=0\).

Take a strictly positive nonconstant polynomial load on each closed
edge, and put \(u=(1+hA_\nu)^{-1}f\). As in Lemma 3.2, \(u\) is analytic
and nonconstant on each edge. It is strictly positive everywhere. An
interior zero contradicts \(-hu''+u=f>0\). At a zero vertex all outgoing
derivatives are nonnegative and sum to zero by (53); each is zero, and
\(u''(p)=-f(p)/h<0\) into each edge again contradicts nonnegativity.
Thus \(m:=\min u>0\).

For a level \(c\in(m,\max u)\) outside the finite set of critical and
vertex values, use \(E_c,\mu,N,a,b\) as in (19). There is at least one
level preimage, by a path from a maximum to a minimum. Integrating over
the superlevel set, with the vertex condition (53), gives

\[
h\left(a(c)+\sum_{u(p)>c}\beta_pu(p)\right)
+U(\mu(c))=\int_{E_c}f\le F(\mu(c)),
\tag{54}
\]

where \(U=\mathcal C_u\) and \(F=\mathcal C_f\). In particular, the
vertex term has a positive sign. Discard it and use \(ab\ge N^2\ge1\)
and \(U''(\mu(c))=-1/b(c)\) to obtain

\[
-hU''+U\le F\quad\text{almost everywhere on }(0,L).
\tag{55}
\]

The rearrangement is Lipschitz by Section 7.1, and only finitely many
ranks are exceptional; thus this is also a weak inequality. The new
boundary information follows by testing the full resolvent equation
against the constant function:

\[
U(L)+h\sum_p\beta_pu(p)=F(L),
\qquad U'(L)=m.
\tag{56}
\]

Since every \(u(p)\ge m\), (56) implies

\[
U(0)=0,\qquad U(L)+h\kappa U'(L)\le F(L).
\tag{57}
\]

Equations (55) and (57), compared to (50) with datum \(f^*\), give
\(U\le V\) by (51)-(52). The polynomial density argument and the
rearrangement contraction (10) pass this inequality to every nonnegative
\(L^2\) load. The last assertion following (52) then gives

\[
\mathcal C_{(1+hA_\nu)^{-1}f}(s)
\le\int_0^s B_{h,\kappa}g(r)\,dr
\quad\text{whenever } f\preccurlyeq g
\tag{58}
\]

for nonnegative \(f\) and nonnegative decreasing \(g\).

\subsection{7.3. General measures and
evolution}\label{general-measures-and-evolution}

Choose finite nonnegative atomic measures \(\nu_j\) converging weakly to
\(\nu\), all with mass \(\kappa\); finite partitions into sets of
diameter tending to zero give such measures. On the fixed graph, the set
of products \(uv\) with \(\|u\|_{H^1},\|v\|_{H^1}\le1\) is relatively
compact in \(C(\Gamma)\). Weak convergence and the uniform bound
\(\|\nu_j-\nu\|_{\mathrm{TV}}\le2\kappa\) therefore imply

\[
\epsilon_j:=\sup_{\|u\|_{H^1},\|v\|_{H^1}\le1}
\left|\int_\Gamma uv\,d(\nu_j-\nu)\right|\longrightarrow0.
\tag{59}
\]

For clarity, take a finite uniform \(\epsilon\)-net of the compact
closure of that product set. The integrals tend to zero on the net, and
the error on the whole set is at most \(2\kappa\epsilon\). Then let
\(\epsilon\) tend to zero.

The bilinear forms \(\langle u,v\rangle+h\mathfrak a_{\nu_j}(u,v)\) are
uniformly coercive on \(H^1\), with coercivity constant at least
\(\min(1,h)\) for the usual norm. Subtracting their weak resolvent
equations, testing with the difference, and applying (59) gives

\[
(1+hA_{\nu_j})^{-1}f\longrightarrow(1+hA_\nu)^{-1}f
\quad\text{in }H^1(\Gamma)
\tag{60}
\]

for every \(f\in L^2(\Gamma)\). For example, if \(u\) is the limiting
resolvent solution, the norm of the difference is bounded by
\(h\epsilon_j\|u\|_{H^1}/\min(1,h)\). Equation (9) now passes (58) to
\(\nu\).

Iterate this resolvent comparison with \(h=t/n\). All interval iterates
stay decreasing and preserve cumulative order by (50)-(52). Spectral
dominated convergence, exactly as in (28), gives (47). Taking \(f=g=1\)
gives (48). This proves the inequality part of Theorem 7.1 without any
regular-level assumption for the heat flow or for a singular measure.

\subsection{7.4. Equality for an arbitrary killing
measure}\label{equality-for-an-arbitrary-killing-measure}

We give the additional details needed because strict edgewise concavity
is generally false in the presence of a distributed potential.

First, every nonnegative \(u\in H^1(\Gamma)\) satisfies the classical
rearrangement inequality

\[
u^*\in H^1(0,L),\qquad
\int_\Gamma|u'|^2\,dx\ge\int_0^L|(u^*)'|^2\,ds.
\tag{61}
\]

One direct extension of the proof in Section 5 is to approximate \(u\)
in \(H^1\) by nonnegative continuous piecewise linear functions on the
edges. At every nonexceptional value their energy coarea formulas and
\(ab\ge N^2\ge1\) prove (61); their finitely many constant pieces
contribute zero energy on both sides. Rearrangement is continuous in
\(L^2\) by (10). The rearranged approximants are bounded in \(H^1\), so
weak compactness and lower semicontinuity pass the estimate to \(u\).
Continuity and connectedness also give \(u^*(L)=\min u\) whenever the
rearrangement has its continuous representative. Hence

\[
\mathfrak a_\nu[u]\ge
\int_0^L|(u^*)'|^2\,ds+\kappa(\min u)^2
=\mathfrak a_{I,\kappa}[u^*].
\tag{62}
\]

Suppose now that equality in (48) holds at \(T>0\). Set
\(u=S_\nu(T/2)1\) and \(z=S_{I,\kappa}(T/2)1\). The half-time identities
(35), with the corresponding forms, and Lemma 5.2 give \(u^*=z\).
Differentiating the nonnegative heat-content deficit at its zero gives

\[
\mathfrak a_\nu[u]=\mathfrak a_{I,\kappa}[z].
\tag{63}
\]

The interval profile \(z\) is strictly positive on \([0,L]\) and
strictly decreasing on \((0,L)\). To see strictness, the
maximum-principle argument in Lemma 5.1 applies to the interval with its
finite Robin endpoint: a zero there would have zero outward derivative
and contradict the boundary point lemma. The semigroup is positivity
improving, is sub-Markovian, and has strictly positive spectrum by (49).
Thus \(A_{I,\kappa}S_{I,\kappa}(t)1>0\) on the open interval as in
(31)-(32). Consequently \(z''<0\) and \(z'(0)=0\), so \(z'(s)<0\) for
\(s>0\). In particular, its distribution has no atoms, including at its
endpoint values.

Equations (62)-(63) force equality separately in the gradient and
killing terms:

\[
\int_\Gamma|u'|^2=\int_0^L|z'|^2,
\qquad
\int_\Gamma\bigl(u^2-m^2\bigr)\,d\nu=0,
\qquad m=\min u=z(L)>0.
\tag{64}
\]

Thus \(\nu\) is supported on the minimum set of \(u\). We next show that
this set is a single endpoint of a path.

Spectral smoothing gives \(u\in\mathcal D(A_\nu)\subset H^1(\Gamma)\).
On each open edge the weak equation is

\[
u''=u\nu-(A_\nu u)\,dx
\tag{65}
\]

as distributions, with the measure restricted to that edge. The right
side is a finite signed measure. Thus \(u'\) has a representative of
bounded variation and is bounded on the edge; \(u\) is edgewise
Lipschitz. This assertion does not assert a sign for \(u''\).

Put \(M=\max u=z(0)\). Since \(u^*=z\), its distribution
\(\mu(c)=|\{u>c\}|\) equals \(z^{-1}(c)\) for \(m<c<M\), and is
continuously differentiable there with \(-\mu'(c)>0\). The
one-dimensional area formula now gives, at almost every such value,
finite level multiplicity \(N(c)\) and

\[
b(c)=\sum_{u=c}|u'|^{-1}=-\mu'(c),
\qquad
\int_\Gamma|u'|^2=\int_m^M a(c)\,dc,
\qquad
\int_0^L|z'|^2=\int_m^M\frac{dc}{b(c)}.
\tag{66}
\]

Here we may exclude vertex values, images of nondifferentiability points
and images of points where \(u'=0\). These images have measure zero:
Lipschitz maps preserve null sets, and the area formula makes the image
of the zero-derivative set null. Moreover that set itself has length
zero in this case. Its pushforward is supported on a null set of levels,
whereas the distribution of \(u\), being that of \(z\), is absolutely
continuous on \((m,M)\) and gives no mass to the endpoint levels. Thus
the reciprocal-derivative formula loses no critical-set mass. Finite
multiplicity almost everywhere follows from
\(\int N(c)dc=\int_\Gamma|u'|dx<\infty\). These facts justify (66)
without a finite-critical-points assertion.

The gradient equality in (64) and \(ab\ge N^2\) yield

\[
0=\int_m^M\left(a(c)-\frac1{b(c)}\right)dc
\ge\int_m^M\frac{N(c)^2-1}{b(c)}\,dc\ge0.
\tag{67}
\]

Hence \(N(c)=1\) for almost every \(c\in(m,M)\). No edge subinterval is
constant, because that would create an atom in the distribution of
\(u\). At a vertex of degree at least three choose disjoint small
neighborhoods of its incidences. On each neighborhood \(u\) takes a
value different from the vertex value. By continuity, each incidence
therefore covers an open interval of values either above or below the
vertex value. Two of the three incidences cover overlapping intervals on
the same side. This produces at least two distinct preimages on an open
interval of levels, contrary to (67). All vertex degrees are at most
two.

A finite connected graph with these degrees is a path or a cycle. On a
cycle, two arcs from a minimum to a maximum give at least two preimages
of every intermediate value, so a cycle is impossible. On a path, a
continuous function with no constant subinterval and almost everywhere
single level multiplicity must be strictly monotone: a failure of
monotonicity produces a strict rise followed by a fall, or a strict fall
followed by a rise, giving two preimages on an open interval of values.
Its minimum is therefore attained at one endpoint only. The support
conclusion from (64) gives \(\nu=\kappa\delta_p\) there. Conversely,
this measure on the path is the comparison problem, proving equality.
Theorem 7.1 follows. \(\square\)

\Needspace{6\baselineskip}\textbf{Corollary 7.2 (integrated comparison).} Under (44),

\[
\int_0^\infty Q_\nu(t)\,dt
\le\frac{L^3}{3}+\frac{L^2}{\kappa},
\tag{68}
\]

with equality exactly in the case of Theorem 7.1. More generally, the
positive-time-weight and convex-temperature conclusions of Section 6
hold with the Robin comparator.

\textbf{Proof.} The positive spectral gap makes the unweighted time
integral finite. The interval torsion function is
\(v(x)=(L^2-x^2)/2+L/\kappa\), whose integral is the right side of (68).
Integrate (48) and use its strictness; the other conclusions follow as
in Section 6. \(\square\)

For atomic vertex measures, (68) and its equality case are precisely
prior work {[}OT, Theorem 8.2{]}, not additional new results. Theorem
7.1 gives their all-times strengthening and allows general measures. The
case \(\kappa=0\) is excluded: every graph then retains the constant
temperature and has heat content \(L\), so rigidity fails. Quantitative
stability and arbitrary vertex couplings remain outside the result. The
two separate extremal comparisons do not settle the product problem
discussed in {[}OT, Section 8.3{]}. Section 8 supplies the required
joint argument for arbitrary finite nonnegative killing measures.

\Needspace{8\baselineskip}\section{8. Kohler-Jobin inequalities for general killing
measures}\label{kohler-jobin-inequalities-for-general-killing-measures}

Let \(\Gamma\) be a finite compact connected metric graph of length
\(L>0\), without Dirichlet conditions, and let \(\nu\) be any finite
nonnegative Borel measure of mass \(\kappa>0\). The form domain is the
continuous space \(H^1(\Gamma)\), as in Section 7. Set

\[
\lambda=\lambda_1(A_\nu),\qquad
T=\int_\Gamma A_\nu^{-1}1\,dx,\qquad s=\kappa L.
\tag{69}
\]

For \(z>0\), let \(\vartheta(z)\in(0,\pi/2)\) solve
\(\vartheta\tan\vartheta=z\). The interval \(I_L\) with Neumann
condition at \(0\) and killing \(\kappa\delta_L\) has eigenvalue
\(\vartheta(s)^2/L^2\) and torsion \(L^3/3+L^2/\kappa\).

\Needspace{6\baselineskip}\textbf{Theorem 8.1 (fixed length and total killing).} Under these
hypotheses,

\[
\lambda T^{2/3}\ge
\vartheta(\kappa L)^2
\left(\frac13+\frac1{\kappa L}\right)^{2/3}.
\tag{70}
\]

Equality holds if and only if \(\Gamma\) is a path and \(\nu\) is
concentrated at one endpoint, up to subdivision and reflection.

\Needspace{6\baselineskip}\textbf{Theorem 8.2 (fixed torsion and total killing).} Let \(\ell>0\)
satisfy \(\ell^3/3+\ell^2/\kappa=T\). Then

\[
\lambda\ge\lambda_1(A_{I_\ell,\kappa}).
\tag{71}
\]

Equivalently, with \(b=\kappa/\sqrt\lambda\) and \(y=\arctan b\),

\[
\lambda^{3/2}T\ge \frac{y^3}{3}+\frac{y^2}{b}.
\tag{72}
\]

The equality case is the same as in Theorem 8.1. Both theorems include
point interactions, integrable nonnegative potentials and singular
continuous killing measures.

These results give sharp fixed-constraint formulations of the
Robin/delta joint-optimization direction discussed in {[}OT, Section 8.3
of the cited preprint{]}, extending the one-point statements of the
preceding revision. The Dirichlet graph Kohler-Jobin theorem is due to
Mugnolo and Plümer {[}MP, Theorem 5.8{]}; its nonlinear Dirichlet
extension is considered in {[}O26{]}. Modified-torsion and level-set
methods are established tools in this subject. Our proof uses a
ground-state distribution functional in which Cauchy-Schwarz absorbs all
killing terms. It then compares a scalar differential inequality with
the cosine profile. Eigenfunction distribution inequalities of this kind
belong to the classical Chiti comparison method; {[}DS, Lemma 3.2{]}
treats Robin domains on the range of interior superlevel sets. Here
connectedness and nonnegative killing give (79) over the full range of
regular levels. The Euclidean domain Robin questions in {[}BCS, Section
5{]} impose different constraints and are not resolved here.

\Needspace{9\baselineskip}\subsection{8.1. A ground-state distribution
estimate}\label{a-ground-state-distribution-estimate}

Define the continuous increasing function

\[
B(x)=\frac{x^3}{3}+x^2\cot x\quad(0<x<\pi/2),
\qquad B(0)=0,\qquad B(\pi/2)=\frac{\pi^3}{24}.
\tag{73}
\]

\Needspace{6\baselineskip}\textbf{Proposition 8.3.} For every graph and measure above, with
\(k=\sqrt\lambda\),

\[
T\ge k^{-3}B\bigl(\min\{kL,\pi/2\}\bigr).
\tag{74}
\]

\textbf{Proof for finitely supported measures.} Subdivide the graph at
the atoms and write \(\nu=\sum_p\beta_p\delta_p\), where
\(\beta_p\ge0\). Choose a positive first eigenfunction \(\psi\), and put
\(m=\min\psi>0\) and \(M=\max\psi\). Positivity follows by taking the
absolute value in the form minimization and applying edgewise
uniqueness: at a zero minimum the nonnegative outgoing derivatives sum
to zero, and propagation along the connected graph would give the zero
function. On each open edge, \(\psi''=-\lambda\psi<0\). Thus no edge
piece is constant, there are finitely many critical and vertex values,
and \(m<M\).

For \(0\le t<M\), define

\[
\mu(t)=|\{\psi>t\}|,\qquad
H(t)=\lambda\int_{\{\psi>t\}}\psi\,dx,
\qquad g(t)=\int_0^t\frac{\mu(r)}{H(r)}\,dr.
\tag{75}
\]

The positive minimum gives \(0\le g'\le1/(\lambda m)\), so \(g\) extends
to a Lipschitz function on \([0,M]\), with \(g(0)=0\). Set
\(G(t)=\int_0^t g'(r)^2\,dr\). Cauchy-Schwarz gives \(g(t)^2\le tG(t)\).
Testing the weak eigenvalue equation with \(G(\psi)\) therefore yields

\[
\begin{aligned}
\mathfrak a_\nu[g(\psi)]
&=\int_\Gamma g'(\psi)^2|\psi'|^2\,dx
  +\int_\Gamma g(\psi)^2\,d\nu\\
&\le\int_\Gamma g'(\psi)^2|\psi'|^2\,dx
  +\int_\Gamma\psi G(\psi)\,d\nu\\
&=\lambda\int_\Gamma\psi G(\psi)\,dx
 =\int_0^M H(t)g'(t)^2\,dt.
\end{aligned}
\tag{76}
\]

Every test function here belongs to \(H^1(\Gamma)\). The chain rules
hold almost everywhere; the finitely many critical points cause no
ambiguity. Cavalieri\textquotesingle s formula and the torsion
variational principle now give

\[
\begin{aligned}
T&=\sup_{v\in H^1(\Gamma)}
 \left(2\int_\Gamma v\,dx-\mathfrak a_\nu[v]\right)\\
&\ge\int_0^M\bigl(2\mu g'-H(g')^2\bigr)\,dt
 =\int_0^M\frac{\mu(t)^2}{H(t)}\,dt.
\end{aligned}
\tag{77}
\]

In particular, killing away from a ground-state minimum is included in
(76); it has not been dropped from the energy.

Let \(q=\psi^*\) be the decreasing rearrangement on \([0,L]\), and put
\(U(s)=\int_0^s q(r)\,dr\). The function \(q\) is positive and
Lipschitz, with \(q(0)=M\) and \(q(L)=m\). To justify Lipschitz
continuity, let \(D_0=\max_\Gamma|\psi'|\). A path crossing two levels
\(t_1<t_2\) spends length at least \((t_2-t_1)/D_0\) between them. Hence
\(\mu(t_1)-\mu(t_2)\ge(t_2-t_1)/D_0\), which gives the asserted inverse
bound. All level and inverse substitutions can also be made on the
finitely many intervals of regular values and then passed to their
endpoints.

\Needspace{5\baselineskip}For a regular \(t\in(m,M)\), flux integration gives

\[
a(t)+\sum_{\psi(p)>t}\beta_p\psi(p)=H(t),
\qquad
 a(t)=\sum_{\psi=t}|\psi'|,\qquad
 b(t)=\sum_{\psi=t}|\psi'|^{-1}=-\mu'(t).
\tag{78}
\]

Connectedness implies that a regular level has at least one preimage, so
\(a(t)b(t)\ge1\). Since the killing coefficients are nonnegative,
\(a(t)\le H(t)\). Using \(H(t)=\lambda U(\mu(t))\) and inverting the
distribution function gives

\[
0\le-q'(s)\le\lambda U(s)\quad\text{for almost every }s\in(0,L).
\tag{79}
\]

The interval \(0<t<m\) in (77) has \(\mu(t)=L\) and
\(H(t)=\lambda U(L)\). Substitution on the remaining levels, followed by
integration by parts, yields

\[
\begin{aligned}
\lambda\int_0^M\frac{\mu(t)^2}{H(t)}\,dt
&=\frac{mL^2}{U(L)}
   +\int_0^L\frac{s^2}{U(s)}(-q'(s))\,ds\\
&=\int_0^L\left(2\frac{s q(s)}{U(s)}
             -\frac{s^2q(s)^2}{U(s)^2}\right)ds.
\end{aligned}
\tag{80}
\]

The term at \(L\) cancels exactly, and \(s^2q(s)/U(s)\to0\) at zero
because \(U(s)/s\to M\). Including the interval below the positive
minimum is essential.

Put \(r(s)=U(s)/q(s)\), with \(r(0)=0\). Since \(q\) is decreasing,
\(r(s)\ge s\). Equation (79) gives

\[
r'(s)=1-\frac{U(s)q'(s)}{q(s)^2}
\le1+k^2r(s)^2,
\qquad
\frac{d}{ds}\arctan(kr(s))\le k.
\tag{81}
\]

Consequently \(r(s)\le\tan(ks)/k\) for \(0<s<\min\{L,\pi/(2k)\}\). Thus
\(z(s)=s/r(s)\) lies in \((0,1]\) everywhere and satisfies
\(z(s)\ge ks\cot(ks)\) on this initial interval. The function \(2z-z^2\)
is nonnegative and increasing for \(0\le z\le1\). Combining (77)-(81)
gives, with \(a=\min\{kL,\pi/2\}\),

\[
\begin{aligned}
T&\ge k^{-2}\int_0^L(2z(s)-z(s)^2)\,ds\\
&\ge k^{-3}\int_0^a(2x\cot x-x^2\cot^2x)\,dx
 = k^{-3}B(a).
\end{aligned}
\tag{82}
\]

The last identity follows by differentiation of (73). This proves (74)
for finite atomic measures.

\textbf{Passage to arbitrary measures.} Choose finite atomic measures
\(\nu_j\rightharpoonup\nu\) of mass \(\kappa\) on the same graph. The
compact embedding \(H^1(\Gamma)\hookrightarrow C(\Gamma)\) and bounded
total masses give

\[
\epsilon_j:=\sup_{\|u\|_{H^1},\|v\|_{H^1}\le1}
\left|\int_\Gamma uv\,d(\nu_j-\nu)\right|\longrightarrow0.
\tag{83}
\]

Indeed, products of two unit-ball functions form a relatively compact
subset of \(C(\Gamma)\), so weak convergence is uniform on that set by a
finite-net argument. The coercivity estimate in Section 7 is uniform in
\(j\). If \(w_j=A_{\nu_j}^{-1}1\) and \(w=A_\nu^{-1}1\), subtracting the
form equations and testing with \(w_j-w\) therefore gives \(w_j\to w\)
in \(H^1\) and \(T_j\to T\).

For completeness, the first eigenvalues converge as well. Testing with a
normalized ground state of \(A_\nu\) gives
\(\limsup\lambda_j\le\lambda\). Conversely, normalized ground states for
\(A_{\nu_j}\) have uniformly bounded \(H^1\) norms; a subsequence
converges weakly in \(H^1\) and strongly in \(C\) and \(L^2\). Its limit
has norm one, and weak lower semicontinuity together with (83) gives
\(\lambda\le\liminf\lambda_j\). Thus \(\lambda_j\to\lambda>0\). The
continuity of the right side of (74) completes the proof. \(\square\)

\Needspace{9\baselineskip}\subsection{8.2. Recovering the prescribed constraints and
equality}\label{recovering-the-prescribed-constraints-and-equality}

We also need the familiar first-eigenvalue comparison

\[
\lambda\ge\lambda_1(A_{I_L,\kappa})
=\frac{\vartheta(\kappa L)^2}{L^2}.
\tag{84}
\]

Here is a direct form proof that fixes the direction. For a nonnegative
\(u\in H^1(\Gamma)\), the graph energy rearrangement proved in Section 7
and the minimum bound give

\[
\mathfrak a_\nu[u]
\ge\int_0^L|(u^*)'|^2\,ds+\kappa u^*(L)^2
\ge\lambda_1(A_{I_L,\kappa})\|u\|_2^2.
\tag{85}
\]

Taking absolute values first handles arbitrary \(u\), and minimization
proves (84).

\Needspace{5\baselineskip}For \(0<x<\pi/2\), differentiation shows

\[
B'(x)=x\cot x\,(2-x\cot x)>0,
\tag{86}
\]

because \(0<x\cot x<1\). For Theorem 8.1, put
\(\theta=\vartheta(\kappa L)\). Equation (84) gives \(kL\ge\theta\).
Therefore (74) and strict monotonicity give

\[
\lambda T^{2/3}
\ge B(\min\{kL,\pi/2\})^{2/3}
\ge B(\theta)^{2/3}
=\theta^2\left(\frac13+\frac1{\kappa L}\right)^{2/3}.
\tag{87}
\]

If equality holds, \(\theta<\pi/2\) and strict monotonicity force
\(kL=\theta\). Equality in (87) then gives \(T=L^3/3+L^2/\kappa\).
Corollary 7.2 characterizes the equality case as the endpoint-killed
interval. Conversely, that interval attains equality.

For Theorem 8.2, let \(b=\kappa/k\) and \(y=\arctan b\). If
\(kL\ge\pi/2\), (74) is strictly stronger than (72). If \(kL<\pi/2\),
(84) and monotonicity of \(x\tan x\) imply

\[
kL\tan(kL)\ge\kappa L,
\qquad kL\ge\arctan(\kappa/k)=y.
\tag{88}
\]

Hence (74) gives \(k^3T\ge B(y)=y^3/3+y^2/b\), proving (72). The
interval of length \(y/k\) with strength \(\kappa\) has first eigenvalue
\(k^2\) and torsion \(k^{-3}B(y)\). Interval torsion is strictly
increasing with length, whereas its first eigenvalue is strictly
decreasing; the latter also follows from
\(k_\ell\ell=\arctan(\kappa/k_\ell)\). Thus (72) is equivalent to (71).
Equality forces \(kL=y\) and again \(T=L^3/3+L^2/\kappa\), giving the
same rigidity by Corollary 7.2. This proves both theorems. \(\square\)

The two general-measure questions left open in revision 1.2 are
therefore settled by (70)-(72). Quantitative stability of these
inequalities and of the heat-content comparison remains open. Section 9
extends the stationary joint comparisons to nonlinear killing. Signed
killing measures, discontinuous vertex couplings, weighted edge forms
and the Euclidean Robin domain problem are outside the present
assertions.

\Needspace{8\baselineskip}\section{\texorpdfstring{9. Nonlinear killing and the
\(p\)-Laplacian}{9. Nonlinear killing and the p-Laplacian}}\label{nonlinear-killing-and-the-p-laplacian}

The distribution argument extends to the full range \(1<p<\infty\). Fix
a finite compact connected metric graph \(\Gamma\) of length \(L\),
without Dirichlet conditions, and a finite nonnegative Borel measure
\(\nu\) with \(\nu(\Gamma)=\kappa>0\). On the continuous space
\(W^{1,p}(\Gamma)\) set

\[
E_{p,\nu}[u]=\int_\Gamma|u'|^p\,dx+\int_\Gamma|u|^p\,d\nu,
\qquad
\lambda_p=\inf_{u\ne0}\frac{E_{p,\nu}[u]}{\int_\Gamma|u|^p\,dx}.
\tag{89}
\]

Let \(w\) minimize \(E_{p,\nu}[u]/p-\int_\Gamma u\,dx\), and distinguish
its mass from the customary homogeneous torsional rigidity:

\[
S_p=\int_\Gamma w\,dx,\qquad T_p=S_p^{p-1},
\qquad p'=\frac p{p-1},\qquad d_p=p'+1.
\tag{90}
\]

This convention for \(T_p\) agrees with {[}O26{]}. In particular, the
scale-invariant product is
\(\lambda_p S_p^{p/d_p}=\lambda_p T_p^{p/(2p-1)}\). The energy includes
nonlinear point interactions as well as distributed killing. The results
below concern stationary variational problems; no nonlinear
heat-concentration theorem is asserted.

We specify the one-dimensional functions to avoid a normalization
ambiguity. Define

\[
A_p(t)=\int_0^t\frac{v^{p-2}}{1+v^p}\,dv,
\qquad a_p=A_p(\infty)=\frac{\pi}{p\sin(\pi/p)},
\qquad \tau_p=A_p^{-1}:[0,a_p)\longrightarrow[0,\infty).
\tag{91}
\]

Thus \(\tau_2=\tan\). For \(\sigma>0\), let
\(\theta_p(\sigma)\in(0,a_p)\) be the unique solution of

\[
\theta_p(\sigma)\tau_p(\theta_p(\sigma))
=\sigma^{1/(p-1)}.
\tag{92}
\]

\Needspace{6\baselineskip}\textbf{Theorem 9.1 (nonlinear joint comparisons).} For every
\(1<p<\infty\) and every graph and measure above, put
\(\sigma=\kappa L^{p-1}\) and \(\theta=\theta_p(\sigma)\). Then

\[
\lambda_p S_p^{p/d_p}
\ge(p-1)\theta^p
\left(\frac1{d_p}+\sigma^{-1/(p-1)}\right)^{p/d_p}.
\tag{93}
\]

Furthermore, if \(\ell>0\) is chosen by

\[
\frac{\ell^{d_p}}{d_p}
+\ell^{p'}\kappa^{-1/(p-1)}=S_p,
\tag{94}
\]

then \(\lambda_p\ge\lambda_p(I_\ell,\kappa\delta_\ell)\). Equality in
either comparison holds exactly when \(\Gamma\) is a path and all
killing is at one endpoint, up to subdivision and reflection.

The nonlinear Dirichlet graph Kohler-Jobin inequality and its
normalization are due to {[}O26, Theorem 5.2{]}, following the nonlinear
modified-torsion strategy of Kohler-Jobin and Brasco. Its level-set flux
inequality and modified-torsion test are established ingredients {[}O26,
Lemma 5.3{]}. Here the Hölder composition estimate below includes all
finite killing terms, and the scalar comparison retains the finite Robin
parameter. Thus the new assertion concerns nonnegative measure killing;
the Dirichlet theorem is credited, not claimed anew.

\Needspace{9\baselineskip}\subsection{9.1. Variational facts and interval
rigidity}\label{variational-facts-and-interval-rigidity}

We first justify the variational objects and the separate interval
comparisons. Path integration and Hölder\textquotesingle s inequality
give, uniformly over measures of mass \(\kappa\),

\[
\|u\|_p^p
\le 2^{p-1}L^p\int_\Gamma|u'|^p\,dx
 +\frac{2^{p-1}L}{\kappa}\int_\Gamma|u|^p\,d\nu.
\tag{95}
\]

Indeed, \(|u(x)|\le |u(y)|+L^{1/p'}\|u'\|_p\); raise to the power \(p\)
and integrate over \(x\) with length measure and over \(y\) with
\(\nu/\kappa\). Consequently \(E_{p,\nu}\) controls the \(W^{1,p}\)
norm. The compact embedding into \(C(\Gamma)\), the direct method and
strict convexity give a unique torsion minimizer \(w\ge0\), a positive
minimum eigenvalue and a nonnegative normalized ground state. Strict
convexity follows because equality in the gradient term makes two
functions differ by a constant, while the positive mass of \(\nu\)
excludes a nonzero constant difference.

The torsion Euler equation and its variational value are

\[
\begin{aligned}
\int_\Gamma|w'|^{p-2}w'v'\,dx
 +\int_\Gamma w^{p-1}v\,d\nu&=\int_\Gamma v\,dx,\\
(p-1)S_p&=\sup_v\left(p\int_\Gamma v\,dx-E_{p,\nu}[v]\right).
\end{aligned}
\tag{96}
\]

Testing with \(w\) gives \(E_{p,\nu}[w]=S_p\). The minimizer is not
zero, since a small positive constant makes \(E_{p,\nu}[u]/p-\int u\)
negative. Optimizing scalar multiples of a test function also gives
\(T_p=\sup_{v\ne0}(\int|v|)^p/E_{p,\nu}[v]\).

For the endpoint-killed interval, direct integration gives

\[
w_I(x)=\frac{L^{p'}-x^{p'}}{p'}
 +\left(\frac L\kappa\right)^{1/(p-1)},
\qquad
S_{p,I}=\frac{L^{d_p}}{d_p}+L^{p'}\kappa^{-1/(p-1)}.
\tag{97}
\]

To compute its eigenvalue, put \(c_p(x)=(1+\tau_p(x)^p)^{-1/p}\).
Differentiation of (91) gives \(c_p'=-\tau_pc_p\) and

\[
-(|c_p'|^{p-2}c_p')'=(p-1)c_p^{p-1},
\qquad c_p(0)=1,\qquad c_p'(0)=0.
\tag{98}
\]

A positive interval ground state decreases from its Neumann endpoint;
its flux equation and initial value therefore identify it, after
normalization, with \(c_p(kx)\), where \(\lambda_p=(p-1)k^p\). Its Robin
condition is \(k^{p-1}\tau_p(kL)^{p-1}=\kappa\), so

\[
\lambda_p(I_L,\kappa\delta_L)
=(p-1)\frac{\theta_p(\kappa L^{p-1})^p}{L^p}.
\tag{99}
\]

\Needspace{6\baselineskip}\textbf{Lemma 9.2 (separate comparisons).} One has
\(\lambda_p\ge\lambda_p(I_L,\kappa\delta_L)\) and \(S_p\le S_{p,I}\).
Equality in the torsion comparison holds exactly for the endpoint-killed
interval.

\textbf{Proof.} The graph rearrangement inequality extends to
\(p\)-energy:

\[
\int_0^L|(u^*)'|^p\,ds\le\int_\Gamma|u'|^p\,dx,
\qquad
E_{p,\nu}[u]\ge E_{p,\kappa\delta_L}[u^*]
\quad(u\ge0).
\tag{100}
\]

For a continuous piecewise affine function, coarea and Hölder give
\(\alpha_pb^{p-1}\ge N^p\ge1\), where
\(\alpha_p(t)=\sum_{u=t}|u'|^{p-1}\) and \(b(t)=\sum_{u=t}|u'|^{-1}\).
Integrating gives the gradient inequality. Piecewise affine
approximation, the \(L^p\) contraction of rearrangement and weak lower
semicontinuity extend it to \(W^{1,p}\), including constant pieces. The
killing term is at least \(\kappa(\min u)^p\), and the continuous
representative of \(u^*\) has endpoint value \(\min u\). This proves
(100). The Rayleigh and torsion variational principles give the two
comparisons.

If \(S_p=S_{p,I}\), uniqueness of the interval torsion minimizer gives
\(w^*=w_I\). Both nonnegative energy defects in (100) vanish, and
\(\nu\) is supported where \(w=\min w\). We explain why the graph is a
path even for singular measures. On each edge the distributional
equation is

\[
(|w'|^{p-2}w')'=w^{p-1}\nu-dx.
\tag{101}
\]

The flux is of bounded variation, hence bounded; therefore \(w'\) is
bounded and \(w\) is Lipschitz. Its rearrangement \(w_I\) is strictly
decreasing, with nonzero derivative on \((0,L]\). Its distribution is
absolutely continuous between endpoint values and has no atoms. The
one-dimensional area formula makes the image of the critical and
nondifferentiability set of \(w\) null. Absolute continuity of the
distribution then shows that this set has zero length. Thus the coarea
energy identities and inverse differentiation apply without losing mass,
as in Section 7.4. Equality in (100) forces \(N(t)=1\) at almost every
level.

There are no constant edge pieces, since they would create atoms in the
distribution. At a vertex with at least three incidences, two disjoint
incident germs have overlapping value intervals on the same side of the
vertex value; those levels have at least two preimages. A cycle likewise
has two arcs joining a minimum to a maximum. Hence the connected graph
has no branch and no cycle, and is a path. A nonmonotone continuous
function on a path would have repeated levels on an interval, so \(w\)
is strictly monotone. Its minimum is one endpoint, where all of \(\nu\)
is concentrated. Formula (97) proves the converse. \(\square\)

\Needspace{9\baselineskip}\subsection{9.2. The distribution bound with nonlinear
killing}\label{the-distribution-bound-with-nonlinear-killing}

Define, continuously at the endpoints,

\[
B_p(x)=\frac{x^{d_p}}{d_p}+\frac{x^{p'}}{\tau_p(x)}
\quad(0<x<a_p),\qquad
B_p(0)=0,\qquad B_p(a_p)=\frac{a_p^{d_p}}{d_p}.
\tag{102}
\]

\Needspace{6\baselineskip}\textbf{Proposition 9.3.} With \(k=(\lambda_p/(p-1))^{1/p}\),

\[
S_p\ge k^{-d_p}B_p(\min\{kL,a_p\}).
\tag{103}
\]

\textbf{Proof for atomic measures.} Subdivide at every atom. A
nonnegative first eigenfunction \(\psi\) is strictly positive. Indeed,
at a zero minimum every outgoing flux is nonnegative and their sum is
zero. On an adjacent edge, the conserved quantity
\((p-1)|\psi'|^p+\lambda_p\psi^p\) is then zero, forcing that edge
solution to vanish; connectedness propagates this contradiction. The
same argument excludes an interior zero. On each edge the flux
derivative is \(-\lambda_p\psi^{p-1}<0\), so the derivative is strictly
decreasing and there is at most one critical point. In particular there
are no constant edge pieces. Write \(m=\min\psi>0\), \(M=\max\psi\), and

\[
\mu(t)=|\{\psi>t\}|,
\qquad H(t)=\lambda_p\int_{\{\psi>t\}}\psi^{p-1}\,dx,
\qquad
 g(t)=\int_0^t\left(\frac{\mu(r)}{H(r)}\right)^{1/(p-1)}dr.
\tag{104}
\]

Since \(H\ge\lambda_p m^{p-1}\mu\), the function \(g\) is Lipschitz on
\([0,M]\). Put \(G(t)=\int_0^t g'(r)^p\,dr\). Hölder gives
\(g(t)^p\le t^{p-1}G(t)\). Testing the weak eigenvalue equation with
\(G(\psi)\), and using the nonnegativity of \(\nu\), gives

\[
\begin{aligned}
E_{p,\nu}[g(\psi)]
&\le\int_\Gamma g'(\psi)^p|\psi'|^p\,dx
       +\int_\Gamma\psi^{p-1}G(\psi)\,d\nu\\
&=\lambda_p\int_\Gamma\psi^{p-1}G(\psi)\,dx
 =\int_0^M H(t)g'(t)^p\,dt.
\end{aligned}
\tag{105}
\]

These compositions are in \(W^{1,p}\) and all chain rules are justified
by their Lipschitz bounds and the finite exceptional set. Equations (96)
and (105), together with Cavalieri\textquotesingle s formula, yield

\[
S_p\ge\int_0^M\frac{\mu(t)^{p'}}{H(t)^{1/(p-1)}}\,dt.
\tag{106}
\]

Indeed the integrand before maximizing is \((p\mu g'-H(g')^p)/(p-1)\),
and (104) attains its pointwise maximum.

Let \(q=\psi^*\) and \(U(s)=\int_0^s q(r)^{p-1}\,dr\). The path argument
of Section 8 makes \(q\) positive and Lipschitz. On regular levels,
total flux, including killing, equals \(H\). Hölder\textquotesingle s
coarea inequality \(\alpha_pb^{p-1}\ge1\) therefore implies

\[
0\le-q'(s)\le(\lambda_p U(s))^{1/(p-1)}
\quad\text{almost everywhere on }(0,L).
\tag{107}
\]

Including the levels \(0<t<m\) in (106), changing to the rank variable
and integrating by parts gives

\[
\begin{aligned}
\lambda_p^{1/(p-1)}
 \int_0^M\frac{\mu(t)^{p'}}{H(t)^{1/(p-1)}}dt
&=\frac{mL^{p'}}{U(L)^{1/(p-1)}}
 +\int_0^L\frac{s^{p'}}{U(s)^{1/(p-1)}}(-q'(s))\,ds\\
&=\int_0^L F_p(z(s))\,ds,\\
z(s)&=\left(\frac{s q(s)^{p-1}}{U(s)}\right)^{1/(p-1)},
\qquad F_p(z)=\frac{pz-z^p}{p-1}.
\end{aligned}
\tag{108}
\]

The term at \(L\) cancels the first term. The term at zero vanishes
because \(U(s)/s\to M^{p-1}\).

Put \(r(s)=U(s)/q(s)^{p-1}\), so \(r\ge s\), \(r(0)=0\) and
\(z=(s/r)^{1/(p-1)}\in(0,1]\). From (107),

\[
r'\le1+(p-1)\lambda_p^{1/(p-1)}r^{p'}
=1+((p-1)kr)^{p'}.
\tag{109}
\]

For \(R=(p-1)kr\), integration of \(R'/(1+R^{p'})\le(p-1)k\) gives
\(\int_0^{R(s)}du/(1+u^{p'})\le(p-1)ks\). Substituting \(u=v^{p-1}\) in
this integral and using (91) shows that, for \(ks<a_p\),

\[
r(s)\le\frac{\tau_p(ks)^{p-1}}{(p-1)k},
\qquad
z(s)\ge\frac{((p-1)ks)^{1/(p-1)}}{\tau_p(ks)}.
\tag{110}
\]

The right side belongs to \((0,1]\): the model ratio \(U/q^{p-1}\) is at
least the rank variable, or equivalently \(\tau_p(x)^{p-1}\ge(p-1)x\),
which follows directly from (91).

The function \(F_p\) is nonnegative and increasing on \([0,1]\).
Differentiating (102), using \(\tau_p'=\tau_p^{2-p}(1+\tau_p^p)\), gives

\[
B_p'(x)=\frac{p'x^{1/(p-1)}}{\tau_p(x)}
 -\frac{x^{p'}}{\tau_p(x)^p}
=(p-1)^{-1/(p-1)}
 F_p\left(\frac{((p-1)x)^{1/(p-1)}}{\tau_p(x)}\right)>0.
\tag{111}
\]

The endpoint values in (102) follow from (91), in particular
\(\tau_p(x)\sim((p-1)x)^{1/(p-1)}\) as \(x\downarrow0\). Equations
(106)-(111), with \(a=\min\{kL,a_p\}\), imply

\[
S_p\ge\lambda_p^{-1/(p-1)}\int_0^{a/k}F_p(z(s))\,ds
\ge k^{-d_p}\int_0^a B_p'(x)\,dx
=k^{-d_p}B_p(a).
\tag{112}
\]

Only the interval before the pole of \(\tau_p\) is compared. The
remaining integrand, when present, is nonnegative. This proves the
atomic case.

\textbf{Passage to arbitrary measures.} Approximate \(\nu\) weakly on
the fixed graph by finite atomic measures \(\nu_j\) of the same mass
\(\kappa\). Compactness of the \(W^{1,p}\) unit ball in \(C(\Gamma)\)
implies uniform convergence of \(\int |u|^p\,d\nu_j\) to
\(\int |u|^p\,d\nu\) on bounded \(W^{1,p}\) sets. The bound (95) is
uniform in \(j\). Testing Rayleigh quotients with a limiting ground
state gives the eigenvalue limsup inequality; compactness of normalized
minimizing ground states, convergence in \(C\) and weak lower
semicontinuity of the gradient energy give the liminf inequality. Thus
\(\lambda_{p,j}\to\lambda_p>0\).

For torsion, the minimizing values of \(E_{p,\nu_j}/p-\int u\) are
uniformly coercive by (95) and Young\textquotesingle s inequality. Every
subsequence of minimizers has a further subsequence converging weakly in
\(W^{1,p}\) and strongly in \(C\). Lower semicontinuity and comparison
with any fixed test function show that the limit minimizes
\(E_{p,\nu}/p-\int u\). Uniqueness identifies it with \(w\), so
\(S_{p,j}\to S_p\). Continuity of the right side of (103), including at
\(kL=a_p\), completes the proof. \(\square\)

\Needspace{9\baselineskip}\subsection{9.3. Proof of the joint
comparisons}\label{proof-of-the-joint-comparisons}

Lemma 9.2 and (99) give \(kL\ge\theta\). Since \(B_p\) is strictly
increasing,

\[
\lambda_p S_p^{p/d_p}
\ge(p-1)B_p(\min\{kL,a_p\})^{p/d_p}
\ge(p-1)B_p(\theta)^{p/d_p}.
\tag{113}
\]

By (92), \(B_p(\theta)=\theta^{d_p}(1/d_p+\sigma^{-1/(p-1)})\), so (113)
is (93). Equality forces \(kL=\theta<a_p\) and \(S_p=S_{p,I}\); Lemma
9.2 gives the stated rigidity.

For the second comparison put \(b=\kappa^{1/(p-1)}/k\) and \(y=A_p(b)\).
If \(kL\ge a_p\), (103) gives \(k^{d_p}S_p>B_p(y)\). Otherwise (99) and
the monotonicity of \(x\tau_p(x)\) imply
\(k\tau_p(kL)\ge\kappa^{1/(p-1)}\), hence \(kL\ge y\). Therefore in all
cases

\[
k^{d_p}S_p\ge B_p(y)
=\frac{y^{d_p}}{d_p}+\frac{y^{p'}}b.
\tag{114}
\]

The interval of length \(y/k\) with killing mass \(\kappa\) has
eigenvalue \((p-1)k^p=\lambda_p\) and torsion mass \(k^{-d_p}B_p(y)\).
At fixed \(\kappa\), (97) strictly increases with length, while the
eigenvalue strictly decreases, as is also seen from
\(k\ell=A_p(\kappa^{1/(p-1)}/k)\). Thus (114) is equivalent to the
comparison with the interval in (94). Equality forces \(kL=y\) and again
\(S_p=S_{p,I}\), proving rigidity. The endpoint interval attains
equality in both statements. \(\square\)

At \(p=2\), all definitions and constants reduce exactly to Section 8.
As \(\sigma\to\infty\), the sharp constant in (93) tends to
\((p-1)a_p^p d_p^{-p/d_p}\), the nonlinear Dirichlet graph constant in
{[}O26, equation (5.4){]}. Quantitative stability, nonlinear heat
concentration and signed-measure analogues remain outside these results.

\raggedright\Needspace{8\baselineskip}\section{References}\label{references}

\Needspace{6\baselineskip}\textbf{{[}DS{]}} Q. Dai and F. Shi, \emph{PPW and Chiti type
inequalities for the eigenvalue problem of Robin Laplacian}, preprint
(2014), \href{https://arxiv.org/abs/1402.2338v1}{arXiv:1402.2338v1}. The
arXiv abstract record uses the variant title \emph{PPW and Chitti type
inequalities for the Robin Laplacian operator}.

\Needspace{6\baselineskip}\textbf{{[}BT{]}} P. Bifulco and M. Täufer, \emph{Faber-Krahn inequality
for the heat content on quantum graphs via random walk expansion},
Electronic Journal of Probability \textbf{30} (2025), paper no. 158,
1-19. \href{https://doi.org/10.1214/25-EJP1414}{DOI:
10.1214/25-EJP1414}. \href{https://arxiv.org/abs/2501.09693}{Preprint
arXiv:2501.09693}. The final Theorem 2.6 and Remark 2.9 specify the
prior result and the all-times question.

\Needspace{6\baselineskip}\textbf{{[}BM{]}} P. Bifulco and D. Mugnolo, \emph{On the Heat Content
of Compact Quantum Graphs}, Annales Henri Poincaré (2026), published
online September 9, 2026.
\href{https://doi.org/10.1007/s00023-026-01755-3}{DOI:
10.1007/s00023-026-01755-3}.
\href{https://arxiv.org/abs/2502.09461}{Preprint arXiv:2502.09461}.

\Needspace{6\baselineskip}\textbf{{[}F{]}} L. Friedlander, \emph{Extremal properties of
eigenvalues for a metric graph}, Annales de l\textquotesingle Institut
Fourier \textbf{55} (2005), no. 1, 199-211.
\href{https://doi.org/10.5802/aif.2095}{DOI: 10.5802/aif.2095}. Lemma 3
and its proof contain the energy rearrangement and level-multiplicity
argument used here.

\Needspace{6\baselineskip}\textbf{{[}V{]}} J. L. Vázquez, \emph{Symmetrization and Mass Comparison
for Degenerate Nonlinear Parabolic and Related Elliptic Equations},
Advanced Nonlinear Studies \textbf{5} (2005), no. 1, 87-131.
\href{https://doi.org/10.1515/ans-2005-0107}{DOI:
10.1515/ans-2005-0107}.
\href{https://verso.mat.uam.es/~juanluis.vazquez/SymmJANS04.pdf}{Author
manuscript, November 25, 2004}, Sections 5.2 and 7, especially Theorems
5.2 and 7.3.

\Needspace{6\baselineskip}\textbf{{[}OT{]}} S. Özcan and M. Täufer, \emph{Torsional Rigidity on
Metric Graphs with Delta-Vertex Conditions}, arXiv:2410.18545v1 (2024).
\href{https://arxiv.org/abs/2410.18545}{Preprint and version history},
especially Theorem 8.2 and Section 8.3. Final publication: Journal of
Mathematical Physics \textbf{67} (2026), 061508,
\href{https://doi.org/10.1063/5.0301900}{DOI: 10.1063/5.0301900}. The
theorem numbering used here is that of the inspected preprint; the final
body was not accessed.

\Needspace{6\baselineskip}\textbf{{[}KS{]}} P. Kurasov and A. Serio, \emph{Optimal Potentials for
Quantum Graphs}, Annales Henri Poincaré \textbf{20} (2019), 1517-1542.
\href{https://doi.org/10.1007/s00023-019-00783-6}{DOI:
10.1007/s00023-019-00783-6}. Cited for the distinct optimization problem
described in its final abstract.

\Needspace{6\baselineskip}\textbf{{[}MP{]}} D. Mugnolo and M. Plümer, \emph{On torsional rigidity
and ground-state energy of compact quantum graphs}, Calculus of
Variations and Partial Differential Equations \textbf{62} (2023),
article 27. \href{https://doi.org/10.1007/s00526-022-02363-9}{DOI:
10.1007/s00526-022-02363-9}. Theorem 5.8 and the modified-torsion
construction in Lemmas 5.15-5.17.

\Needspace{6\baselineskip}\textbf{{[}BCS{]}} G. Buttazzo, S. Cito and F. Solombrino,
\emph{Relations Between Principal Eigenvalue and Torsional Rigidity with
Robin Boundary Conditions}, Milan Journal of Mathematics \textbf{94}
(2026), 369-386. \href{https://doi.org/10.1007/s00032-026-00438-2}{DOI:
10.1007/s00032-026-00438-2}. See Section 5 for domain product questions
with fixed Robin parameter and volume.

\Needspace{6\baselineskip}\textbf{{[}O26{]}} S. Özcan, \emph{On the p-torsional rigidity of
compact metric graphs: a sharp Kohler-Jobin inequality},
arXiv:2607.12333v2 (July 15, 2026).
\href{https://arxiv.org/abs/2607.12333}{Preprint and version history}.
Theorem 5.2 gives the nonlinear Dirichlet graph inequality; Lemmas
5.3-5.4 develop its modified-torsion proof. The paper assumes a nonempty
Dirichlet set and nonlinear Kirchhoff conditions.

\Needspace{8\baselineskip}\section{Research transparency}\label{research-transparency}

This is an AI-assisted research preprint prepared under Henry
Zweiman\textquotesingle s name at his request. The proof and the
primary-source comparison were audited internally by the originating
assistant; no independent expert review or journal acceptance is
claimed. The accompanying assessment and source records state the search
scope, the source passages inspected and the remaining priority
limitations. The mathematical statements above are supported by the
displayed arguments, rather than by finite computations or artifact
checks.

\end{document}
